% A short paper exercising the LaTeX constructs the extractor has to separate
% from prose. Every misspelling is deliberate; expected.md lists them.

\documentclass[11pt]{article}
\usepackage[utf8]{inputenc}
\usepackage{amsmath}
\usepackage{algorithm}
\usepackage{listings}
\usepackage{hyperref}

\title{On the Separation of Markup from Prose}
\author{A. Reader}

\begin{document}
\maketitle

\begin{abstract}
  An abstract is prose and is checked. This sentence contains a deliberatly
  wrong word so that the example has something to report. The title and the
  author above are not checked: the whole preamble is skipped, so that a name
  in \verb|\author{}| is not reported as a misspelling every time.
\end{abstract}

\section{Introduction}
\label{sec:intro}

Body text is prose. A citation such as \cite{knuth1984} is not, and neither is
the label above or the reference in Section~\ref{sec:intro}. A command whose
argument is prose, such as \emph{emphasised text}, does not break the sentence
around it.

Inline mathematics, $a^2 + b^2 = c^2$, sits inside a sentence without splitting
it. Display mathematics does not:
\begin{equation}
  \int_0^1 f(x)\,\mathrm{d}x = \sum_{i=1}^{n} a_i
  \label{eq:sum}
\end{equation}

\subsection{Environments that hold no prose}

% skip_environments in .languagecheck.yaml adds `algorithm` to the built-in set,
% so nothing inside it is checked -- including the misspelling below.
\begin{algorithm}
  Compute the anser by iterating until convergence.
\end{algorithm}

\begin{lstlisting}[language=Python]
# A comment in a listing: never checked.
recieve = "misspelt on purpose"
\end{lstlisting}

\subsection{Turning the checker off}

% lang-check-disable-next-line
This line holds a deliberate mispelling and is not reported.

% lang-check-begin spelling.typo
Inside this region only spelling is silenced; every other rule still applies.
Here is one more mispelling.
% lang-check-end

% A `lang:` directive routes a passage to another language. LaTeX has no
% built-in language declaration the way Typst does, so the directive is how a
% quotation in another language is marked.
% lang-check-begin lang:fr
Voici un paragraphe en français au milieu d'un article anglais, avec une
fautte volontaire.
% lang-check-end

\section{Conclusion}

A final sentence with a deliberatly wrong word, and a URL that is not prose:
\url{https://example.org/paper}.

\begin{thebibliography}{1}
\bibitem{knuth1984}
  Donald E. Knuth. \emph{The \TeX book}. Addison-Wesley, 1984.
\end{thebibliography}

\end{document}
